\documentclass[11pt]{article}
\usepackage[margin=1in]{geometry}
\usepackage{graphicx}
\usepackage{hyperref}
\usepackage{booktabs}

\title{AMD-Agent: Continual Malware Detection}
\author{}
\date{\today}

\begin{document}
\maketitle

\begin{abstract}
LangGraph-orchestrated continual learning for Windows PE malware: autonomous collection, static feature extraction, River ADWIN concept drift monitoring, LightGBM classification, MADAR replay retraining, and TESSERACT temporal evaluation with drift evidence logging.
\end{abstract}

\section{Introduction}
Malware detectors degrade as attack patterns shift---a classic concept drift problem. Static PE classifiers trained on a fixed snapshot cannot adapt without new labeled data and controlled retraining. This work presents AMD-Agent, a continual learning pipeline that autonomously ingests PE samples from threat feeds and APIs, monitors distributional shift, and retrains when drift is detected. The system is implemented as a LangGraph \texttt{StateGraph} with explicit acquisition, analysis, and learning phases, culminating in an \texttt{evaluation} node that records metrics on a configurable cadence and after retraining.

\section{Architecture}
Figure~\ref{fig:flow} summarizes the LangGraph flow. Both terminal branches (\texttt{classifier\_inference} and \texttt{model\_retrain}) converge on \texttt{evaluation} before \texttt{END}.

\begin{figure}[h]
\centering
\fbox{\parbox{0.9\linewidth}{
\small
START $\rightarrow$ SourceSelector $\rightarrow$ (ThreatIntelIngest $|$ SourceDiscovery) $\rightarrow$ BinaryFetch $\rightarrow$ DataValidation $\rightarrow$ FeatureExtraction $\rightarrow$ DriftMonitor\\
$\rightarrow$ [no drift] ClassifierInference $\rightarrow$ Evaluation $\rightarrow$ END\\
$\rightarrow$ [drift] ExplainDriftContext $\rightarrow$ ModelRetrain $\rightarrow$ Evaluation $\rightarrow$ END
}}
\caption{LangGraph state flow (steady phase). Bootstrap skips threat-intel ingest and drift monitoring.}
\label{fig:flow}
\end{figure}

Shared state is a Pydantic \texttt{AgentState} (hashes, features, drift flags, \texttt{evaluation\_metrics}, \texttt{drift\_pre\_metrics}, \texttt{pending\_drift\_log}).

\section{Data Provenance and Autonomy Boundaries}
Table~\ref{tab:provenance} separates dynamic collection from fixed configuration. Course experiments use \texttt{AMD\_ALLOW\_LOCAL\_BENIGN=0} and an empty \texttt{data/benign/} directory.

\begin{table}[h]
\centering
\caption{Data provenance}
\label{tab:provenance}
\begin{tabular}{@{}ll@{}}
\toprule
\textbf{Dynamic (runtime)} & \textbf{Hardcoded / curated} \\
\midrule
MalwareBazaar / ThreatFox / OTX pulse API discovery & Bootstrap malware provider pin (MalwareBazaar) \\
OTX pulse hash evidence + structured LLM verdicts & Curated native CTI seed list \\
Curated CTI hash extraction from allowlisted feeds & \texttt{GITHUB\_BENIGN\_REPOS} allowlist \\
Sysinternals directory scrape & CTI download domain allowlist \\
SQLite tracker labels from provider metadata & Optional \texttt{data/benign/} if env enabled \\
\bottomrule
\end{tabular}
\end{table}

\section{Methodology}
\textbf{Bootstrap.} Repeated graph passes until $\geq 100$ malware and $\geq 100$ benign trainable samples exist; drift monitoring is disabled.

\textbf{Steady state.} Malware flows through threat-intel ingest or provider discovery; benign via Sysinternals/GitHub. Static features (2304-dim EMBER-like vector) are stored in SQLite.

\textbf{Drift.} River ADWIN on mean section entropy plus rolling multivariate shift on selected model features. On detection, the latest persisted TESSERACT metrics are attached as pre-drift context; after MADAR retrain, post-metrics are written to \texttt{drift\_log.jsonl}.

\textbf{Classification.} LightGBM with Optuna tuning and FPR-aware threshold ($\leq 0.1\%$ target).

\textbf{Evaluation.} Chronological train/val/test splits (no random holdout); metrics appended to \texttt{evaluation\_log.jsonl}; AUT over accuracy history. Bootstrap training uses a stratified split for stability, while temporal TESSERACT remains the evaluation view.

\section{Results}
\IfFileExists{figures/performance_decay.png}{%
\begin{figure}[h]
\centering
\includegraphics[width=0.8\textwidth]{figures/performance_decay.png}
\caption{TESSERACT temporal performance decay (accuracy and FPR vs.\ evaluation run).}
\end{figure}
}{%
Performance decay plot: run the daemon, then copy \texttt{data/figures/performance\_decay.png} to \texttt{report/figures/performance\_decay.png}.%
}

Drift cycles are documented in \texttt{drift\_log.jsonl} (pre/post accuracy and FPR, \texttt{mean\_shift}, \texttt{corr\_shift}, retrain flag). Area Under Time (AUT) summarizes recovery across retraining iterations when multiple evaluation runs are available.

\section{Limitations}
\begin{itemize}
  \item Requires MalwareBazaar API key; Ollama and capa are optional but improve source selection and drift explanations.
  \item Drift detection is inactive during bootstrap; sparse labels can skip verified drift batches.
  \item TESSERACT is skipped when temporal splits are single-class or sample counts are too low.
  \item Steady-state malware still depends heavily on curated intel feeds and MalwareBazaar availability.
\end{itemize}

\section{Conclusion}
AMD-Agent demonstrates a Docker-first continual malware pipeline with explicit drift-triggered retraining and first-class evaluation logging suitable for temporal analysis. Future work includes LangGraph subgraphs, broader bootstrap provider rotation, and tighter integration of LLM tool calls with ingest actions.

\end{document}
